\chapter{Introduction}
\label{chap:introduction}

\section{Background and Motivation}

Online learning studies sequential decision problems in which a learner
repeatedly chooses among available actions without knowing future rewards or
losses. A standard objective is to minimize \emph{external regret}, which
compares the learner with the best fixed action in hindsight. Classical work
includes Hannan's repeated-play framework \citep{dresher_approximation_1958},
the Weighted Majority algorithm \citep{littlestone_weighted_1994}, and Hedge
\citep{freund_decision-theoretic_1997}. Under bandit feedback, where only the
reward of the selected action is observed, algorithms such as EXP3 provide the
corresponding adversarial guarantee \citep{auer_nonstochastic_2002}.

External regret considers only fixed deviations. Internal regret allows one
action to be replaced by another, while swap regret allows a separate
replacement for every action. These stronger notions are particularly relevant
in repeated games: vanishing external regret is associated with coarse
correlated equilibrium \citep{chen_hedging_2020}, whereas vanishing internal
or swap regret is associated with correlated equilibrium
\citep{hart_simple_2000,blum_external_2007}. This has motivated both direct
regret-matching procedures and reductions that construct stronger-regret
learners from external-regret algorithms
\citep{hart_simple_2000,blum_external_2007,ito_tight_2020}.

Theoretical guarantees, however, do not by themselves determine finite-horizon
behavior. Algorithms with the same asymptotic order can differ substantially
over practical horizons, and stronger guarantees may come with different
finite-horizon constants or update structures. The environment also matters:
an adaptive reward process can react to the learner's past actions, whereas an
oblivious process is fixed independently of them. In repeated games, the
opponent is itself part of the reward-generation process.

Structured self-play can introduce additional regularity beyond an arbitrary
adversarial sequence. For example, optimistic methods can obtain faster regret
rates in repeated games \citep{chen_hedging_2020}, and recent work shows that
no-swap-regret dynamics can have stronger convergence properties in symmetric
zero-sum self-play \citep{leme_convergence_2024}. These results provide useful
theoretical context, but they do not directly answer how the different
algorithms compare empirically under a common finite-horizon evaluator.

This thesis therefore focuses on an empirical comparison. The goal is to study
how external-, internal-, and swap-regret algorithms behave under the same
experimental framework, how stable the observed differences are across
feedback models and reward processes, and how they translate into repeated-game
behavior.

\section{Research Questions and Contributions}

The empirical study is organized around three research questions:

\begin{enumerate}
    \item \textbf{How do the considered algorithms differ in external,
    internal, and swap regret at finite horizons?}

    \item \textbf{How stable are these differences across full-information and
    bandit feedback, adaptive and oblivious reward processes, action-space
    sizes, and horizons?}

    \item \textbf{In repeated games, how are the regret differences reflected
    in CE/CCE behavior, and how strongly do they depend on the opponent?}
\end{enumerate}

To address these questions, the thesis implements and evaluates a collection of
existing external-regret, regret-matching, and swap-regret methods under a
common framework. The experiments use expected action regret as the common
regret metric. They cover full-information and bandit feedback, an adaptive
historical-frequency environment, an oblivious lazy-random-walk environment,
action-space and horizon scaling, and self-play and cross-play in RPS and
RPSLS. The repeated-game experiments additionally evaluate the empirical
joint-action distribution relative to the correlated-equilibrium and
coarse-correlated-equilibrium sets.

The thesis does not propose a new learning algorithm or derive new regret
bounds. Its contribution is the consistent empirical comparison of existing
methods and the analysis of their finite-horizon behavior across these
settings.

\section{Thesis Structure}

Chapter~2 introduces the online-learning setting, the regret notions used in
the thesis, and the equilibrium concepts needed for the repeated-game
experiments. Chapter~3 presents the implemented algorithms and the literature
results most relevant to their interpretation. Chapter~4 describes the
experimental environments, configuration, and evaluation metrics.
Chapter~5 reports the empirical results. Chapter~6 discusses the main
finite-horizon findings and relates them to the existing literature. Finally,
Chapter~7 summarizes the conclusions and outlines directions for future work.
